\chapter{Scattering, Phase Space, and Cross Sections}

\chapterquote{A field-theory formula is useful only after the smaller calculation inside it is visible.}

This chapter is part of \emph{Interactions and Perturbative Quantum Field Theory}.  The working vocabulary is S-matrix, rates, phase space, normalization.  The first pass through the chapter should be slow: identify the object being changed, the condition held fixed, and the reason a boundary term or normalization is allowed.  The first concrete theme is expanding the interaction exponential; later sections reuse the same habit in a more compact notation.

For Scattering, Phase Space, and Cross Sections, the calculus-level starting point is tied to expanding the interaction exponential.  Matrices, waves, operators, fields, and gauge variables are introduced only as the calculation requires them.  A formula is accepted after it has been checked in a finite model, differentiated or varied explicitly, and connected to the field-theoretic use that motivates it.

\section{The concrete problem: expanding the interaction exponential}

The work of this section is expanding the interaction exponential.  In the chapter heading the vocabulary is S-matrix, rates, phase space, normalization, but the first objects on the page are more prosaic: interaction terms, contractions, propagators, vertices, diagrams, external legs, and amplitudes.  The formal notation will be useful only after it has been tied to a limit, a symmetry, a normalization condition, or an integration-by-parts step.  In Scattering, Phase Space, and Cross Sections, section 1, the useful habit is to name the object, the operation, and the assumption before using the compact notation.

The finite model is a Gaussian average corrected by a small non-quadratic term expanded as a power series.  In that model no mystery is available: every derivative is an ordinary derivative with respect to one listed variable, every inner product is a sum, and every boundary assumption can be written at the endpoints.  This is why the finite model is not a toy in the dismissive sense.  It is the bookkeeping device that prevents continuum notation from becoming a fog of indices.  For Scattering, Phase Space, and Cross Sections, section 1, that bookkeeping is deliberately modest, but it is what later keeps signs, factors of \(2\pi\), and normalization constants from turning into guesses.

The central idea for this chapter is each source derivative chooses a field, each contraction gives a propagator, and each interaction term gives a vertex.  To test that idea, strip the formula down until only the operation remains.  A sign change after raising or lowering an index means the metric has entered.  A derivative moving from one factor to another means a boundary term has been handled.  A normalization constant means that some unit object, such as a delta function, basis vector, or one-particle state, has been fixed.  Read in the setting of Scattering, Phase Space, and Cross Sections, section 1, the line is a check on the calculation rather than an invitation to memorize another rule.

For the concrete problem: expanding the interaction exponential, the calculation should also pass a dimensional check.  The left and right sides must carry the same units; more subtly, they must transform in the same way under the symmetries already introduced.  This is not a proof, but it is a quick way to catch wrong factors of \(2\pi\), signs in the metric, missing powers of the lattice spacing, or an omitted charge.  The same step will look more abstract later; in Scattering, Phase Space, and Cross Sections, section 1 it is kept close to the finite calculation so the limiting move remains visible.

The bridge to quantum field theory is direct.  Feynman rules, scattering amplitudes, LSZ reduction, and loop corrections are organized perturbative expansions of correlation functions.  Later notation will carry more indices and fewer verbal reminders, so the safe habit is to keep asking which operation is being performed and which quantity is being held fixed.  At this point in Scattering, Phase Space, and Cross Sections, section 1, it is worth asking what would fail if the boundary condition, normalization, or symmetry assumption were changed.

One warning is worth making explicit here.  A diagram is not the calculation; it is a bookkeeping device for a definite integral and symmetry factor.  This warning points to the delicate step of the derivation, not to a decorative aside.  In Scattering, Phase Space, and Cross Sections, section 1, the calculation is meant to leave a trail from an ordinary derivative, sum, matrix product, or Gaussian integral to the field-theory formula.

The section closes by keeping expanding the interaction exponential connected to Scattering, Phase Space, and Cross Sections.  The same general operation will be used with different objects in neighboring chapters, so the present calculation is both a result and a rehearsal.  In Scattering, Phase Space, and Cross Sections, section 1, the useful habit is to name the object, the operation, and the assumption before using the compact notation.



\paragraph{Step-by-step derivation.}

For Scattering, Phase Space, and Cross Sections: The concrete problem: expanding the interaction exponential, the derivation is written from the smallest usable model upward.

Split the action into a free quadratic part and an interaction.  In the path integral,
\[
  e^{\ii S}=e^{\ii S_0}\left(1+\ii S_I+\frac{\ii^2}{2}S_I^2+\cdots\right).
\]
For \(\Lag_I=-\lambda\phi^4/4!\), one insertion contributes \(-\ii\lambda\int\dd^4x\,\phi(x)^4/4!\).  The four fields at that point can contract with four external fields in \(4!\) ways, cancelling the \(4!\) in the Lagrangian.  The vertex factor is therefore \(-\ii\lambda\).  In Scattering, Phase Space, and Cross Sections, section 1, the useful habit is to name the object, the operation, and the assumption before using the compact notation.

External propagators are removed by LSZ because asymptotic particles are identified with the residues of poles in correlation functions.  What remains after the pole factors are stripped is the scattering amplitude \(\mathcal M\).  For Scattering, Phase Space, and Cross Sections, section 1, that bookkeeping is deliberately modest, but it is what later keeps signs, factors of \(2\pi\), and normalization constants from turning into guesses.

The formulas to keep in view are

\[
  \Delta_F(p)=\frac{\ii}{p^2-m^2+\ii\epsilon}
\]
\[
  e^{\ii\int\Lag_I}=1+\ii\int\Lag_I+\frac{\ii^2}{2}\left(\int\Lag_I\right)^2+\cdots
\]
\[
  \ii\mathcal M=\hbox{sum of connected amputated diagrams}
\]

In this setting the calculation for Scattering, Phase Space, and Cross Sections: The concrete problem: expanding the interaction exponential is complete when the finite model, the limiting step, and the normalization or boundary condition can all be named without looking back at the prose.




\begin{example}[A worked calculation for Scattering, Phase Space, and Cross Sections: The concrete problem: expanding the interaction exponential]
In \(\lambda\phi^4\) theory, the four-point function at first order has one interaction insertion.  The four fields at the vertex can be paired with four external fields in \(4!\) ways, cancelling the \(1/4!\).  The remaining factor is \(-\ii\lambda\), multiplied by four propagators from the vertex point to the external points.  In Scattering, Phase Space, and Cross Sections, section 1, the useful habit is to name the object, the operation, and the assumption before using the compact notation.
\end{example}



\begin{exercise}
Why does the factor \(1/4!\) in \(-\lambda\phi^4/4!\) disappear from the four-point vertex rule?  Use the notation of Scattering, Phase Space, and Cross Sections: The concrete problem: expanding the interaction exponential.
\begin{answer}
There are \(4!\) contractions connecting the four identical fields at the vertex to four external fields.
\end{answer}
\end{exercise}

\begin{exercise}
State what LSZ removes from a connected correlation function.  Use the notation of Scattering, Phase Space, and Cross Sections: The concrete problem: expanding the interaction exponential.
\begin{answer}
It removes the external propagator poles and leaves the on-shell scattering amplitude.
\end{answer}
\end{exercise}



\section{Objects, notation, and units: proving Wick contractions}

The work of this section is proving Wick contractions.  In the chapter heading the vocabulary is S-matrix, rates, phase space, normalization, but the first objects on the page are more prosaic: interaction terms, contractions, propagators, vertices, diagrams, external legs, and amplitudes.  The formal notation will be useful only after it has been tied to a limit, a symmetry, a normalization condition, or an integration-by-parts step.  In Scattering, Phase Space, and Cross Sections, section 2, the useful habit is to name the object, the operation, and the assumption before using the compact notation.

The finite model is a Gaussian average corrected by a small non-quadratic term expanded as a power series.  In that model no mystery is available: every derivative is an ordinary derivative with respect to one listed variable, every inner product is a sum, and every boundary assumption can be written at the endpoints.  This is why the finite model is not a toy in the dismissive sense.  It is the bookkeeping device that prevents continuum notation from becoming a fog of indices.  For Scattering, Phase Space, and Cross Sections, section 2, that bookkeeping is deliberately modest, but it is what later keeps signs, factors of \(2\pi\), and normalization constants from turning into guesses.

The central idea for this chapter is each source derivative chooses a field, each contraction gives a propagator, and each interaction term gives a vertex.  To test that idea, strip the formula down until only the operation remains.  A sign change after raising or lowering an index means the metric has entered.  A derivative moving from one factor to another means a boundary term has been handled.  A normalization constant means that some unit object, such as a delta function, basis vector, or one-particle state, has been fixed.  Read in the setting of Scattering, Phase Space, and Cross Sections, section 2, the line is a check on the calculation rather than an invitation to memorize another rule.

For objects, notation, and units: proving wick contractions, the calculation should also pass a dimensional check.  The left and right sides must carry the same units; more subtly, they must transform in the same way under the symmetries already introduced.  This is not a proof, but it is a quick way to catch wrong factors of \(2\pi\), signs in the metric, missing powers of the lattice spacing, or an omitted charge.  The same step will look more abstract later; in Scattering, Phase Space, and Cross Sections, section 2 it is kept close to the finite calculation so the limiting move remains visible.

The bridge to quantum field theory is direct.  Feynman rules, scattering amplitudes, LSZ reduction, and loop corrections are organized perturbative expansions of correlation functions.  Later notation will carry more indices and fewer verbal reminders, so the safe habit is to keep asking which operation is being performed and which quantity is being held fixed.  At this point in Scattering, Phase Space, and Cross Sections, section 2, it is worth asking what would fail if the boundary condition, normalization, or symmetry assumption were changed.

One warning is worth making explicit here.  A diagram is not the calculation; it is a bookkeeping device for a definite integral and symmetry factor.  This warning points to the delicate step of the derivation, not to a decorative aside.  In Scattering, Phase Space, and Cross Sections, section 2, the calculation is meant to leave a trail from an ordinary derivative, sum, matrix product, or Gaussian integral to the field-theory formula.

The section closes by keeping proving Wick contractions connected to Scattering, Phase Space, and Cross Sections.  The same general operation will be used with different objects in neighboring chapters, so the present calculation is both a result and a rehearsal.  In Scattering, Phase Space, and Cross Sections, section 2, the useful habit is to name the object, the operation, and the assumption before using the compact notation.



\paragraph{Step-by-step derivation.}

For Scattering, Phase Space, and Cross Sections: Objects, notation, and units: proving Wick contractions, the derivation is written from the smallest usable model upward.

Split the action into a free quadratic part and an interaction.  In the path integral,
\[
  e^{\ii S}=e^{\ii S_0}\left(1+\ii S_I+\frac{\ii^2}{2}S_I^2+\cdots\right).
\]
For \(\Lag_I=-\lambda\phi^4/4!\), one insertion contributes \(-\ii\lambda\int\dd^4x\,\phi(x)^4/4!\).  The four fields at that point can contract with four external fields in \(4!\) ways, cancelling the \(4!\) in the Lagrangian.  The vertex factor is therefore \(-\ii\lambda\).  In Scattering, Phase Space, and Cross Sections, section 2, the useful habit is to name the object, the operation, and the assumption before using the compact notation.

External propagators are removed by LSZ because asymptotic particles are identified with the residues of poles in correlation functions.  What remains after the pole factors are stripped is the scattering amplitude \(\mathcal M\).  For Scattering, Phase Space, and Cross Sections, section 2, that bookkeeping is deliberately modest, but it is what later keeps signs, factors of \(2\pi\), and normalization constants from turning into guesses.

The formulas to keep in view are

\[
  \Delta_F(p)=\frac{\ii}{p^2-m^2+\ii\epsilon}
\]
\[
  e^{\ii\int\Lag_I}=1+\ii\int\Lag_I+\frac{\ii^2}{2}\left(\int\Lag_I\right)^2+\cdots
\]
\[
  \ii\mathcal M=\hbox{sum of connected amputated diagrams}
\]

In this setting the calculation for Scattering, Phase Space, and Cross Sections: Objects, notation, and units: proving Wick contractions is complete when the finite model, the limiting step, and the normalization or boundary condition can all be named without looking back at the prose.




\begin{example}[A worked calculation for Scattering, Phase Space, and Cross Sections: Objects, notation, and units: proving Wick contractions]
In \(\lambda\phi^4\) theory, the four-point function at first order has one interaction insertion.  The four fields at the vertex can be paired with four external fields in \(4!\) ways, cancelling the \(1/4!\).  The remaining factor is \(-\ii\lambda\), multiplied by four propagators from the vertex point to the external points.  In Scattering, Phase Space, and Cross Sections, section 2, the useful habit is to name the object, the operation, and the assumption before using the compact notation.
\end{example}



\begin{exercise}
Why does the factor \(1/4!\) in \(-\lambda\phi^4/4!\) disappear from the four-point vertex rule?  Use the notation of Scattering, Phase Space, and Cross Sections: Objects, notation, and units: proving Wick contractions.
\begin{answer}
There are \(4!\) contractions connecting the four identical fields at the vertex to four external fields.
\end{answer}
\end{exercise}

\begin{exercise}
State what LSZ removes from a connected correlation function.  Use the notation of Scattering, Phase Space, and Cross Sections: Objects, notation, and units: proving Wick contractions.
\begin{answer}
It removes the external propagator poles and leaves the on-shell scattering amplitude.
\end{answer}
\end{exercise}



\section{The finite calculation: reading a vertex factor}

The work of this section is reading a vertex factor.  In the chapter heading the vocabulary is S-matrix, rates, phase space, normalization, but the first objects on the page are more prosaic: interaction terms, contractions, propagators, vertices, diagrams, external legs, and amplitudes.  The formal notation will be useful only after it has been tied to a limit, a symmetry, a normalization condition, or an integration-by-parts step.  In Scattering, Phase Space, and Cross Sections, section 3, the useful habit is to name the object, the operation, and the assumption before using the compact notation.

The finite model is a Gaussian average corrected by a small non-quadratic term expanded as a power series.  In that model no mystery is available: every derivative is an ordinary derivative with respect to one listed variable, every inner product is a sum, and every boundary assumption can be written at the endpoints.  This is why the finite model is not a toy in the dismissive sense.  It is the bookkeeping device that prevents continuum notation from becoming a fog of indices.  For Scattering, Phase Space, and Cross Sections, section 3, that bookkeeping is deliberately modest, but it is what later keeps signs, factors of \(2\pi\), and normalization constants from turning into guesses.

The central idea for this chapter is each source derivative chooses a field, each contraction gives a propagator, and each interaction term gives a vertex.  To test that idea, strip the formula down until only the operation remains.  A sign change after raising or lowering an index means the metric has entered.  A derivative moving from one factor to another means a boundary term has been handled.  A normalization constant means that some unit object, such as a delta function, basis vector, or one-particle state, has been fixed.  Read in the setting of Scattering, Phase Space, and Cross Sections, section 3, the line is a check on the calculation rather than an invitation to memorize another rule.

For the finite calculation: reading a vertex factor, the calculation should also pass a dimensional check.  The left and right sides must carry the same units; more subtly, they must transform in the same way under the symmetries already introduced.  This is not a proof, but it is a quick way to catch wrong factors of \(2\pi\), signs in the metric, missing powers of the lattice spacing, or an omitted charge.  The same step will look more abstract later; in Scattering, Phase Space, and Cross Sections, section 3 it is kept close to the finite calculation so the limiting move remains visible.

The bridge to quantum field theory is direct.  Feynman rules, scattering amplitudes, LSZ reduction, and loop corrections are organized perturbative expansions of correlation functions.  Later notation will carry more indices and fewer verbal reminders, so the safe habit is to keep asking which operation is being performed and which quantity is being held fixed.  At this point in Scattering, Phase Space, and Cross Sections, section 3, it is worth asking what would fail if the boundary condition, normalization, or symmetry assumption were changed.

One warning is worth making explicit here.  A diagram is not the calculation; it is a bookkeeping device for a definite integral and symmetry factor.  This warning points to the delicate step of the derivation, not to a decorative aside.  In Scattering, Phase Space, and Cross Sections, section 3, the calculation is meant to leave a trail from an ordinary derivative, sum, matrix product, or Gaussian integral to the field-theory formula.

The section closes by keeping reading a vertex factor connected to Scattering, Phase Space, and Cross Sections.  The same general operation will be used with different objects in neighboring chapters, so the present calculation is both a result and a rehearsal.  In Scattering, Phase Space, and Cross Sections, section 3, the useful habit is to name the object, the operation, and the assumption before using the compact notation.



\paragraph{Step-by-step derivation.}

For Scattering, Phase Space, and Cross Sections: The finite calculation: reading a vertex factor, the derivation is written from the smallest usable model upward.

Split the action into a free quadratic part and an interaction.  In the path integral,
\[
  e^{\ii S}=e^{\ii S_0}\left(1+\ii S_I+\frac{\ii^2}{2}S_I^2+\cdots\right).
\]
For \(\Lag_I=-\lambda\phi^4/4!\), one insertion contributes \(-\ii\lambda\int\dd^4x\,\phi(x)^4/4!\).  The four fields at that point can contract with four external fields in \(4!\) ways, cancelling the \(4!\) in the Lagrangian.  The vertex factor is therefore \(-\ii\lambda\).  In Scattering, Phase Space, and Cross Sections, section 3, the useful habit is to name the object, the operation, and the assumption before using the compact notation.

External propagators are removed by LSZ because asymptotic particles are identified with the residues of poles in correlation functions.  What remains after the pole factors are stripped is the scattering amplitude \(\mathcal M\).  For Scattering, Phase Space, and Cross Sections, section 3, that bookkeeping is deliberately modest, but it is what later keeps signs, factors of \(2\pi\), and normalization constants from turning into guesses.

The formulas to keep in view are

\[
  \Delta_F(p)=\frac{\ii}{p^2-m^2+\ii\epsilon}
\]
\[
  e^{\ii\int\Lag_I}=1+\ii\int\Lag_I+\frac{\ii^2}{2}\left(\int\Lag_I\right)^2+\cdots
\]
\[
  \ii\mathcal M=\hbox{sum of connected amputated diagrams}
\]

In this setting the calculation for Scattering, Phase Space, and Cross Sections: The finite calculation: reading a vertex factor is complete when the finite model, the limiting step, and the normalization or boundary condition can all be named without looking back at the prose.




\begin{example}[A worked calculation for Scattering, Phase Space, and Cross Sections: The finite calculation: reading a vertex factor]
In \(\lambda\phi^4\) theory, the four-point function at first order has one interaction insertion.  The four fields at the vertex can be paired with four external fields in \(4!\) ways, cancelling the \(1/4!\).  The remaining factor is \(-\ii\lambda\), multiplied by four propagators from the vertex point to the external points.  In Scattering, Phase Space, and Cross Sections, section 3, the useful habit is to name the object, the operation, and the assumption before using the compact notation.
\end{example}



\begin{exercise}
Why does the factor \(1/4!\) in \(-\lambda\phi^4/4!\) disappear from the four-point vertex rule?  Use the notation of Scattering, Phase Space, and Cross Sections: The finite calculation: reading a vertex factor.
\begin{answer}
There are \(4!\) contractions connecting the four identical fields at the vertex to four external fields.
\end{answer}
\end{exercise}

\begin{exercise}
State what LSZ removes from a connected correlation function.  Use the notation of Scattering, Phase Space, and Cross Sections: The finite calculation: reading a vertex factor.
\begin{answer}
It removes the external propagator poles and leaves the on-shell scattering amplitude.
\end{answer}
\end{exercise}



\section{The central derivation: conserving momentum at a vertex}

The work of this section is conserving momentum at a vertex.  In the chapter heading the vocabulary is S-matrix, rates, phase space, normalization, but the first objects on the page are more prosaic: interaction terms, contractions, propagators, vertices, diagrams, external legs, and amplitudes.  The formal notation will be useful only after it has been tied to a limit, a symmetry, a normalization condition, or an integration-by-parts step.  In Scattering, Phase Space, and Cross Sections, section 4, the useful habit is to name the object, the operation, and the assumption before using the compact notation.

The finite model is a Gaussian average corrected by a small non-quadratic term expanded as a power series.  In that model no mystery is available: every derivative is an ordinary derivative with respect to one listed variable, every inner product is a sum, and every boundary assumption can be written at the endpoints.  This is why the finite model is not a toy in the dismissive sense.  It is the bookkeeping device that prevents continuum notation from becoming a fog of indices.  For Scattering, Phase Space, and Cross Sections, section 4, that bookkeeping is deliberately modest, but it is what later keeps signs, factors of \(2\pi\), and normalization constants from turning into guesses.

The central idea for this chapter is each source derivative chooses a field, each contraction gives a propagator, and each interaction term gives a vertex.  To test that idea, strip the formula down until only the operation remains.  A sign change after raising or lowering an index means the metric has entered.  A derivative moving from one factor to another means a boundary term has been handled.  A normalization constant means that some unit object, such as a delta function, basis vector, or one-particle state, has been fixed.  Read in the setting of Scattering, Phase Space, and Cross Sections, section 4, the line is a check on the calculation rather than an invitation to memorize another rule.

For the central derivation: conserving momentum at a vertex, the calculation should also pass a dimensional check.  The left and right sides must carry the same units; more subtly, they must transform in the same way under the symmetries already introduced.  This is not a proof, but it is a quick way to catch wrong factors of \(2\pi\), signs in the metric, missing powers of the lattice spacing, or an omitted charge.  The same step will look more abstract later; in Scattering, Phase Space, and Cross Sections, section 4 it is kept close to the finite calculation so the limiting move remains visible.

The bridge to quantum field theory is direct.  Feynman rules, scattering amplitudes, LSZ reduction, and loop corrections are organized perturbative expansions of correlation functions.  Later notation will carry more indices and fewer verbal reminders, so the safe habit is to keep asking which operation is being performed and which quantity is being held fixed.  At this point in Scattering, Phase Space, and Cross Sections, section 4, it is worth asking what would fail if the boundary condition, normalization, or symmetry assumption were changed.

One warning is worth making explicit here.  A diagram is not the calculation; it is a bookkeeping device for a definite integral and symmetry factor.  This warning points to the delicate step of the derivation, not to a decorative aside.  In Scattering, Phase Space, and Cross Sections, section 4, the calculation is meant to leave a trail from an ordinary derivative, sum, matrix product, or Gaussian integral to the field-theory formula.

The section closes by keeping conserving momentum at a vertex connected to Scattering, Phase Space, and Cross Sections.  The same general operation will be used with different objects in neighboring chapters, so the present calculation is both a result and a rehearsal.  In Scattering, Phase Space, and Cross Sections, section 4, the useful habit is to name the object, the operation, and the assumption before using the compact notation.



\paragraph{Step-by-step derivation.}

For Scattering, Phase Space, and Cross Sections: The central derivation: conserving momentum at a vertex, the derivation is written from the smallest usable model upward.

Split the action into a free quadratic part and an interaction.  In the path integral,
\[
  e^{\ii S}=e^{\ii S_0}\left(1+\ii S_I+\frac{\ii^2}{2}S_I^2+\cdots\right).
\]
For \(\Lag_I=-\lambda\phi^4/4!\), one insertion contributes \(-\ii\lambda\int\dd^4x\,\phi(x)^4/4!\).  The four fields at that point can contract with four external fields in \(4!\) ways, cancelling the \(4!\) in the Lagrangian.  The vertex factor is therefore \(-\ii\lambda\).  In Scattering, Phase Space, and Cross Sections, section 4, the useful habit is to name the object, the operation, and the assumption before using the compact notation.

External propagators are removed by LSZ because asymptotic particles are identified with the residues of poles in correlation functions.  What remains after the pole factors are stripped is the scattering amplitude \(\mathcal M\).  For Scattering, Phase Space, and Cross Sections, section 4, that bookkeeping is deliberately modest, but it is what later keeps signs, factors of \(2\pi\), and normalization constants from turning into guesses.

The formulas to keep in view are

\[
  \Delta_F(p)=\frac{\ii}{p^2-m^2+\ii\epsilon}
\]
\[
  e^{\ii\int\Lag_I}=1+\ii\int\Lag_I+\frac{\ii^2}{2}\left(\int\Lag_I\right)^2+\cdots
\]
\[
  \ii\mathcal M=\hbox{sum of connected amputated diagrams}
\]

In this setting the calculation for Scattering, Phase Space, and Cross Sections: The central derivation: conserving momentum at a vertex is complete when the finite model, the limiting step, and the normalization or boundary condition can all be named without looking back at the prose.




\begin{example}[A worked calculation for Scattering, Phase Space, and Cross Sections: The central derivation: conserving momentum at a vertex]
In \(\lambda\phi^4\) theory, the four-point function at first order has one interaction insertion.  The four fields at the vertex can be paired with four external fields in \(4!\) ways, cancelling the \(1/4!\).  The remaining factor is \(-\ii\lambda\), multiplied by four propagators from the vertex point to the external points.  In Scattering, Phase Space, and Cross Sections, section 4, the useful habit is to name the object, the operation, and the assumption before using the compact notation.
\end{example}



\begin{exercise}
Why does the factor \(1/4!\) in \(-\lambda\phi^4/4!\) disappear from the four-point vertex rule?  Use the notation of Scattering, Phase Space, and Cross Sections: The central derivation: conserving momentum at a vertex.
\begin{answer}
There are \(4!\) contractions connecting the four identical fields at the vertex to four external fields.
\end{answer}
\end{exercise}

\begin{exercise}
State what LSZ removes from a connected correlation function.  Use the notation of Scattering, Phase Space, and Cross Sections: The central derivation: conserving momentum at a vertex.
\begin{answer}
It removes the external propagator poles and leaves the on-shell scattering amplitude.
\end{answer}
\end{exercise}



\section{Worked calculation: amputating external propagators}

The work of this section is amputating external propagators.  In the chapter heading the vocabulary is S-matrix, rates, phase space, normalization, but the first objects on the page are more prosaic: interaction terms, contractions, propagators, vertices, diagrams, external legs, and amplitudes.  The formal notation will be useful only after it has been tied to a limit, a symmetry, a normalization condition, or an integration-by-parts step.  In Scattering, Phase Space, and Cross Sections, section 5, the useful habit is to name the object, the operation, and the assumption before using the compact notation.

The finite model is a Gaussian average corrected by a small non-quadratic term expanded as a power series.  In that model no mystery is available: every derivative is an ordinary derivative with respect to one listed variable, every inner product is a sum, and every boundary assumption can be written at the endpoints.  This is why the finite model is not a toy in the dismissive sense.  It is the bookkeeping device that prevents continuum notation from becoming a fog of indices.  For Scattering, Phase Space, and Cross Sections, section 5, that bookkeeping is deliberately modest, but it is what later keeps signs, factors of \(2\pi\), and normalization constants from turning into guesses.

The central idea for this chapter is each source derivative chooses a field, each contraction gives a propagator, and each interaction term gives a vertex.  To test that idea, strip the formula down until only the operation remains.  A sign change after raising or lowering an index means the metric has entered.  A derivative moving from one factor to another means a boundary term has been handled.  A normalization constant means that some unit object, such as a delta function, basis vector, or one-particle state, has been fixed.  Read in the setting of Scattering, Phase Space, and Cross Sections, section 5, the line is a check on the calculation rather than an invitation to memorize another rule.

For worked calculation: amputating external propagators, the calculation should also pass a dimensional check.  The left and right sides must carry the same units; more subtly, they must transform in the same way under the symmetries already introduced.  This is not a proof, but it is a quick way to catch wrong factors of \(2\pi\), signs in the metric, missing powers of the lattice spacing, or an omitted charge.  The same step will look more abstract later; in Scattering, Phase Space, and Cross Sections, section 5 it is kept close to the finite calculation so the limiting move remains visible.

The bridge to quantum field theory is direct.  Feynman rules, scattering amplitudes, LSZ reduction, and loop corrections are organized perturbative expansions of correlation functions.  Later notation will carry more indices and fewer verbal reminders, so the safe habit is to keep asking which operation is being performed and which quantity is being held fixed.  At this point in Scattering, Phase Space, and Cross Sections, section 5, it is worth asking what would fail if the boundary condition, normalization, or symmetry assumption were changed.

One warning is worth making explicit here.  A diagram is not the calculation; it is a bookkeeping device for a definite integral and symmetry factor.  This warning points to the delicate step of the derivation, not to a decorative aside.  In Scattering, Phase Space, and Cross Sections, section 5, the calculation is meant to leave a trail from an ordinary derivative, sum, matrix product, or Gaussian integral to the field-theory formula.

The section closes by keeping amputating external propagators connected to Scattering, Phase Space, and Cross Sections.  The same general operation will be used with different objects in neighboring chapters, so the present calculation is both a result and a rehearsal.  In Scattering, Phase Space, and Cross Sections, section 5, the useful habit is to name the object, the operation, and the assumption before using the compact notation.



\paragraph{Step-by-step derivation.}

For Scattering, Phase Space, and Cross Sections: Worked calculation: amputating external propagators, the derivation is written from the smallest usable model upward.

Split the action into a free quadratic part and an interaction.  In the path integral,
\[
  e^{\ii S}=e^{\ii S_0}\left(1+\ii S_I+\frac{\ii^2}{2}S_I^2+\cdots\right).
\]
For \(\Lag_I=-\lambda\phi^4/4!\), one insertion contributes \(-\ii\lambda\int\dd^4x\,\phi(x)^4/4!\).  The four fields at that point can contract with four external fields in \(4!\) ways, cancelling the \(4!\) in the Lagrangian.  The vertex factor is therefore \(-\ii\lambda\).  In Scattering, Phase Space, and Cross Sections, section 5, the useful habit is to name the object, the operation, and the assumption before using the compact notation.

External propagators are removed by LSZ because asymptotic particles are identified with the residues of poles in correlation functions.  What remains after the pole factors are stripped is the scattering amplitude \(\mathcal M\).  For Scattering, Phase Space, and Cross Sections, section 5, that bookkeeping is deliberately modest, but it is what later keeps signs, factors of \(2\pi\), and normalization constants from turning into guesses.

The formulas to keep in view are

\[
  \Delta_F(p)=\frac{\ii}{p^2-m^2+\ii\epsilon}
\]
\[
  e^{\ii\int\Lag_I}=1+\ii\int\Lag_I+\frac{\ii^2}{2}\left(\int\Lag_I\right)^2+\cdots
\]
\[
  \ii\mathcal M=\hbox{sum of connected amputated diagrams}
\]

In this setting the calculation for Scattering, Phase Space, and Cross Sections: Worked calculation: amputating external propagators is complete when the finite model, the limiting step, and the normalization or boundary condition can all be named without looking back at the prose.




\begin{example}[A worked calculation for Scattering, Phase Space, and Cross Sections: Worked calculation: amputating external propagators]
In \(\lambda\phi^4\) theory, the four-point function at first order has one interaction insertion.  The four fields at the vertex can be paired with four external fields in \(4!\) ways, cancelling the \(1/4!\).  The remaining factor is \(-\ii\lambda\), multiplied by four propagators from the vertex point to the external points.  In Scattering, Phase Space, and Cross Sections, section 5, the useful habit is to name the object, the operation, and the assumption before using the compact notation.
\end{example}



\begin{exercise}
Why does the factor \(1/4!\) in \(-\lambda\phi^4/4!\) disappear from the four-point vertex rule?  Use the notation of Scattering, Phase Space, and Cross Sections: Worked calculation: amputating external propagators.
\begin{answer}
There are \(4!\) contractions connecting the four identical fields at the vertex to four external fields.
\end{answer}
\end{exercise}

\begin{exercise}
State what LSZ removes from a connected correlation function.  Use the notation of Scattering, Phase Space, and Cross Sections: Worked calculation: amputating external propagators.
\begin{answer}
It removes the external propagator poles and leaves the on-shell scattering amplitude.
\end{answer}
\end{exercise}



\section{Boundary conditions, normalization, and signs: normalizing phase space}

The work of this section is normalizing phase space.  In the chapter heading the vocabulary is S-matrix, rates, phase space, normalization, but the first objects on the page are more prosaic: interaction terms, contractions, propagators, vertices, diagrams, external legs, and amplitudes.  The formal notation will be useful only after it has been tied to a limit, a symmetry, a normalization condition, or an integration-by-parts step.  In Scattering, Phase Space, and Cross Sections, section 6, the useful habit is to name the object, the operation, and the assumption before using the compact notation.

The finite model is a Gaussian average corrected by a small non-quadratic term expanded as a power series.  In that model no mystery is available: every derivative is an ordinary derivative with respect to one listed variable, every inner product is a sum, and every boundary assumption can be written at the endpoints.  This is why the finite model is not a toy in the dismissive sense.  It is the bookkeeping device that prevents continuum notation from becoming a fog of indices.  For Scattering, Phase Space, and Cross Sections, section 6, that bookkeeping is deliberately modest, but it is what later keeps signs, factors of \(2\pi\), and normalization constants from turning into guesses.

The central idea for this chapter is each source derivative chooses a field, each contraction gives a propagator, and each interaction term gives a vertex.  To test that idea, strip the formula down until only the operation remains.  A sign change after raising or lowering an index means the metric has entered.  A derivative moving from one factor to another means a boundary term has been handled.  A normalization constant means that some unit object, such as a delta function, basis vector, or one-particle state, has been fixed.  Read in the setting of Scattering, Phase Space, and Cross Sections, section 6, the line is a check on the calculation rather than an invitation to memorize another rule.

For boundary conditions, normalization, and signs: normalizing phase space, the calculation should also pass a dimensional check.  The left and right sides must carry the same units; more subtly, they must transform in the same way under the symmetries already introduced.  This is not a proof, but it is a quick way to catch wrong factors of \(2\pi\), signs in the metric, missing powers of the lattice spacing, or an omitted charge.  The same step will look more abstract later; in Scattering, Phase Space, and Cross Sections, section 6 it is kept close to the finite calculation so the limiting move remains visible.

The bridge to quantum field theory is direct.  Feynman rules, scattering amplitudes, LSZ reduction, and loop corrections are organized perturbative expansions of correlation functions.  Later notation will carry more indices and fewer verbal reminders, so the safe habit is to keep asking which operation is being performed and which quantity is being held fixed.  At this point in Scattering, Phase Space, and Cross Sections, section 6, it is worth asking what would fail if the boundary condition, normalization, or symmetry assumption were changed.

One warning is worth making explicit here.  A diagram is not the calculation; it is a bookkeeping device for a definite integral and symmetry factor.  This warning points to the delicate step of the derivation, not to a decorative aside.  In Scattering, Phase Space, and Cross Sections, section 6, the calculation is meant to leave a trail from an ordinary derivative, sum, matrix product, or Gaussian integral to the field-theory formula.

The section closes by keeping normalizing phase space connected to Scattering, Phase Space, and Cross Sections.  The same general operation will be used with different objects in neighboring chapters, so the present calculation is both a result and a rehearsal.  In Scattering, Phase Space, and Cross Sections, section 6, the useful habit is to name the object, the operation, and the assumption before using the compact notation.



\paragraph{Step-by-step derivation.}

For Scattering, Phase Space, and Cross Sections: Boundary conditions, normalization, and signs: normalizing phase space, the derivation is written from the smallest usable model upward.

Split the action into a free quadratic part and an interaction.  In the path integral,
\[
  e^{\ii S}=e^{\ii S_0}\left(1+\ii S_I+\frac{\ii^2}{2}S_I^2+\cdots\right).
\]
For \(\Lag_I=-\lambda\phi^4/4!\), one insertion contributes \(-\ii\lambda\int\dd^4x\,\phi(x)^4/4!\).  The four fields at that point can contract with four external fields in \(4!\) ways, cancelling the \(4!\) in the Lagrangian.  The vertex factor is therefore \(-\ii\lambda\).  In Scattering, Phase Space, and Cross Sections, section 6, the useful habit is to name the object, the operation, and the assumption before using the compact notation.

External propagators are removed by LSZ because asymptotic particles are identified with the residues of poles in correlation functions.  What remains after the pole factors are stripped is the scattering amplitude \(\mathcal M\).  For Scattering, Phase Space, and Cross Sections, section 6, that bookkeeping is deliberately modest, but it is what later keeps signs, factors of \(2\pi\), and normalization constants from turning into guesses.

The formulas to keep in view are

\[
  \Delta_F(p)=\frac{\ii}{p^2-m^2+\ii\epsilon}
\]
\[
  e^{\ii\int\Lag_I}=1+\ii\int\Lag_I+\frac{\ii^2}{2}\left(\int\Lag_I\right)^2+\cdots
\]
\[
  \ii\mathcal M=\hbox{sum of connected amputated diagrams}
\]

In this setting the calculation for Scattering, Phase Space, and Cross Sections: Boundary conditions, normalization, and signs: normalizing phase space is complete when the finite model, the limiting step, and the normalization or boundary condition can all be named without looking back at the prose.




\begin{example}[A worked calculation for Scattering, Phase Space, and Cross Sections: Boundary conditions, normalization, and signs: normalizing phase space]
In \(\lambda\phi^4\) theory, the four-point function at first order has one interaction insertion.  The four fields at the vertex can be paired with four external fields in \(4!\) ways, cancelling the \(1/4!\).  The remaining factor is \(-\ii\lambda\), multiplied by four propagators from the vertex point to the external points.  In Scattering, Phase Space, and Cross Sections, section 6, the useful habit is to name the object, the operation, and the assumption before using the compact notation.
\end{example}



\begin{exercise}
Why does the factor \(1/4!\) in \(-\lambda\phi^4/4!\) disappear from the four-point vertex rule?  Use the notation of Scattering, Phase Space, and Cross Sections: Boundary conditions, normalization, and signs: normalizing phase space.
\begin{answer}
There are \(4!\) contractions connecting the four identical fields at the vertex to four external fields.
\end{answer}
\end{exercise}

\begin{exercise}
State what LSZ removes from a connected correlation function.  Use the notation of Scattering, Phase Space, and Cross Sections: Boundary conditions, normalization, and signs: normalizing phase space.
\begin{answer}
It removes the external propagator poles and leaves the on-shell scattering amplitude.
\end{answer}
\end{exercise}



\section{How the idea enters quantum field theory: identifying loop variables}

The work of this section is identifying loop variables.  In the chapter heading the vocabulary is S-matrix, rates, phase space, normalization, but the first objects on the page are more prosaic: interaction terms, contractions, propagators, vertices, diagrams, external legs, and amplitudes.  The formal notation will be useful only after it has been tied to a limit, a symmetry, a normalization condition, or an integration-by-parts step.  In Scattering, Phase Space, and Cross Sections, section 7, the useful habit is to name the object, the operation, and the assumption before using the compact notation.

The finite model is a Gaussian average corrected by a small non-quadratic term expanded as a power series.  In that model no mystery is available: every derivative is an ordinary derivative with respect to one listed variable, every inner product is a sum, and every boundary assumption can be written at the endpoints.  This is why the finite model is not a toy in the dismissive sense.  It is the bookkeeping device that prevents continuum notation from becoming a fog of indices.  For Scattering, Phase Space, and Cross Sections, section 7, that bookkeeping is deliberately modest, but it is what later keeps signs, factors of \(2\pi\), and normalization constants from turning into guesses.

The central idea for this chapter is each source derivative chooses a field, each contraction gives a propagator, and each interaction term gives a vertex.  To test that idea, strip the formula down until only the operation remains.  A sign change after raising or lowering an index means the metric has entered.  A derivative moving from one factor to another means a boundary term has been handled.  A normalization constant means that some unit object, such as a delta function, basis vector, or one-particle state, has been fixed.  Read in the setting of Scattering, Phase Space, and Cross Sections, section 7, the line is a check on the calculation rather than an invitation to memorize another rule.

For how the idea enters quantum field theory: identifying loop variables, the calculation should also pass a dimensional check.  The left and right sides must carry the same units; more subtly, they must transform in the same way under the symmetries already introduced.  This is not a proof, but it is a quick way to catch wrong factors of \(2\pi\), signs in the metric, missing powers of the lattice spacing, or an omitted charge.  The same step will look more abstract later; in Scattering, Phase Space, and Cross Sections, section 7 it is kept close to the finite calculation so the limiting move remains visible.

The bridge to quantum field theory is direct.  Feynman rules, scattering amplitudes, LSZ reduction, and loop corrections are organized perturbative expansions of correlation functions.  Later notation will carry more indices and fewer verbal reminders, so the safe habit is to keep asking which operation is being performed and which quantity is being held fixed.  At this point in Scattering, Phase Space, and Cross Sections, section 7, it is worth asking what would fail if the boundary condition, normalization, or symmetry assumption were changed.

One warning is worth making explicit here.  A diagram is not the calculation; it is a bookkeeping device for a definite integral and symmetry factor.  This warning points to the delicate step of the derivation, not to a decorative aside.  In Scattering, Phase Space, and Cross Sections, section 7, the calculation is meant to leave a trail from an ordinary derivative, sum, matrix product, or Gaussian integral to the field-theory formula.

The section closes by keeping identifying loop variables connected to Scattering, Phase Space, and Cross Sections.  The same general operation will be used with different objects in neighboring chapters, so the present calculation is both a result and a rehearsal.  In Scattering, Phase Space, and Cross Sections, section 7, the useful habit is to name the object, the operation, and the assumption before using the compact notation.



\paragraph{Step-by-step derivation.}

For Scattering, Phase Space, and Cross Sections: How the idea enters quantum field theory: identifying loop variables, the derivation is written from the smallest usable model upward.

Split the action into a free quadratic part and an interaction.  In the path integral,
\[
  e^{\ii S}=e^{\ii S_0}\left(1+\ii S_I+\frac{\ii^2}{2}S_I^2+\cdots\right).
\]
For \(\Lag_I=-\lambda\phi^4/4!\), one insertion contributes \(-\ii\lambda\int\dd^4x\,\phi(x)^4/4!\).  The four fields at that point can contract with four external fields in \(4!\) ways, cancelling the \(4!\) in the Lagrangian.  The vertex factor is therefore \(-\ii\lambda\).  In Scattering, Phase Space, and Cross Sections, section 7, the useful habit is to name the object, the operation, and the assumption before using the compact notation.

External propagators are removed by LSZ because asymptotic particles are identified with the residues of poles in correlation functions.  What remains after the pole factors are stripped is the scattering amplitude \(\mathcal M\).  For Scattering, Phase Space, and Cross Sections, section 7, that bookkeeping is deliberately modest, but it is what later keeps signs, factors of \(2\pi\), and normalization constants from turning into guesses.

The formulas to keep in view are

\[
  \Delta_F(p)=\frac{\ii}{p^2-m^2+\ii\epsilon}
\]
\[
  e^{\ii\int\Lag_I}=1+\ii\int\Lag_I+\frac{\ii^2}{2}\left(\int\Lag_I\right)^2+\cdots
\]
\[
  \ii\mathcal M=\hbox{sum of connected amputated diagrams}
\]

In this setting the calculation for Scattering, Phase Space, and Cross Sections: How the idea enters quantum field theory: identifying loop variables is complete when the finite model, the limiting step, and the normalization or boundary condition can all be named without looking back at the prose.




\begin{example}[A worked calculation for Scattering, Phase Space, and Cross Sections: How the idea enters quantum field theory: identifying loop variables]
In \(\lambda\phi^4\) theory, the four-point function at first order has one interaction insertion.  The four fields at the vertex can be paired with four external fields in \(4!\) ways, cancelling the \(1/4!\).  The remaining factor is \(-\ii\lambda\), multiplied by four propagators from the vertex point to the external points.  In Scattering, Phase Space, and Cross Sections, section 7, the useful habit is to name the object, the operation, and the assumption before using the compact notation.
\end{example}



\begin{exercise}
Why does the factor \(1/4!\) in \(-\lambda\phi^4/4!\) disappear from the four-point vertex rule?  Use the notation of Scattering, Phase Space, and Cross Sections: How the idea enters quantum field theory: identifying loop variables.
\begin{answer}
There are \(4!\) contractions connecting the four identical fields at the vertex to four external fields.
\end{answer}
\end{exercise}

\begin{exercise}
State what LSZ removes from a connected correlation function.  Use the notation of Scattering, Phase Space, and Cross Sections: How the idea enters quantum field theory: identifying loop variables.
\begin{answer}
It removes the external propagator poles and leaves the on-shell scattering amplitude.
\end{answer}
\end{exercise}



\section{Review problems and extensions: checking symmetry factors}

The work of this section is checking symmetry factors.  In the chapter heading the vocabulary is S-matrix, rates, phase space, normalization, but the first objects on the page are more prosaic: interaction terms, contractions, propagators, vertices, diagrams, external legs, and amplitudes.  The formal notation will be useful only after it has been tied to a limit, a symmetry, a normalization condition, or an integration-by-parts step.  In Scattering, Phase Space, and Cross Sections, section 8, the useful habit is to name the object, the operation, and the assumption before using the compact notation.

The finite model is a Gaussian average corrected by a small non-quadratic term expanded as a power series.  In that model no mystery is available: every derivative is an ordinary derivative with respect to one listed variable, every inner product is a sum, and every boundary assumption can be written at the endpoints.  This is why the finite model is not a toy in the dismissive sense.  It is the bookkeeping device that prevents continuum notation from becoming a fog of indices.  For Scattering, Phase Space, and Cross Sections, section 8, that bookkeeping is deliberately modest, but it is what later keeps signs, factors of \(2\pi\), and normalization constants from turning into guesses.

The central idea for this chapter is each source derivative chooses a field, each contraction gives a propagator, and each interaction term gives a vertex.  To test that idea, strip the formula down until only the operation remains.  A sign change after raising or lowering an index means the metric has entered.  A derivative moving from one factor to another means a boundary term has been handled.  A normalization constant means that some unit object, such as a delta function, basis vector, or one-particle state, has been fixed.  Read in the setting of Scattering, Phase Space, and Cross Sections, section 8, the line is a check on the calculation rather than an invitation to memorize another rule.

For review problems and extensions: checking symmetry factors, the calculation should also pass a dimensional check.  The left and right sides must carry the same units; more subtly, they must transform in the same way under the symmetries already introduced.  This is not a proof, but it is a quick way to catch wrong factors of \(2\pi\), signs in the metric, missing powers of the lattice spacing, or an omitted charge.  The same step will look more abstract later; in Scattering, Phase Space, and Cross Sections, section 8 it is kept close to the finite calculation so the limiting move remains visible.

The bridge to quantum field theory is direct.  Feynman rules, scattering amplitudes, LSZ reduction, and loop corrections are organized perturbative expansions of correlation functions.  Later notation will carry more indices and fewer verbal reminders, so the safe habit is to keep asking which operation is being performed and which quantity is being held fixed.  At this point in Scattering, Phase Space, and Cross Sections, section 8, it is worth asking what would fail if the boundary condition, normalization, or symmetry assumption were changed.

One warning is worth making explicit here.  A diagram is not the calculation; it is a bookkeeping device for a definite integral and symmetry factor.  This warning points to the delicate step of the derivation, not to a decorative aside.  In Scattering, Phase Space, and Cross Sections, section 8, the calculation is meant to leave a trail from an ordinary derivative, sum, matrix product, or Gaussian integral to the field-theory formula.

The section closes by keeping checking symmetry factors connected to Scattering, Phase Space, and Cross Sections.  The same general operation will be used with different objects in neighboring chapters, so the present calculation is both a result and a rehearsal.  In Scattering, Phase Space, and Cross Sections, section 8, the useful habit is to name the object, the operation, and the assumption before using the compact notation.



\paragraph{Step-by-step derivation.}

For Scattering, Phase Space, and Cross Sections: Review problems and extensions: checking symmetry factors, the derivation is written from the smallest usable model upward.

Split the action into a free quadratic part and an interaction.  In the path integral,
\[
  e^{\ii S}=e^{\ii S_0}\left(1+\ii S_I+\frac{\ii^2}{2}S_I^2+\cdots\right).
\]
For \(\Lag_I=-\lambda\phi^4/4!\), one insertion contributes \(-\ii\lambda\int\dd^4x\,\phi(x)^4/4!\).  The four fields at that point can contract with four external fields in \(4!\) ways, cancelling the \(4!\) in the Lagrangian.  The vertex factor is therefore \(-\ii\lambda\).  In Scattering, Phase Space, and Cross Sections, section 8, the useful habit is to name the object, the operation, and the assumption before using the compact notation.

External propagators are removed by LSZ because asymptotic particles are identified with the residues of poles in correlation functions.  What remains after the pole factors are stripped is the scattering amplitude \(\mathcal M\).  For Scattering, Phase Space, and Cross Sections, section 8, that bookkeeping is deliberately modest, but it is what later keeps signs, factors of \(2\pi\), and normalization constants from turning into guesses.

The formulas to keep in view are

\[
  \Delta_F(p)=\frac{\ii}{p^2-m^2+\ii\epsilon}
\]
\[
  e^{\ii\int\Lag_I}=1+\ii\int\Lag_I+\frac{\ii^2}{2}\left(\int\Lag_I\right)^2+\cdots
\]
\[
  \ii\mathcal M=\hbox{sum of connected amputated diagrams}
\]

In this setting the calculation for Scattering, Phase Space, and Cross Sections: Review problems and extensions: checking symmetry factors is complete when the finite model, the limiting step, and the normalization or boundary condition can all be named without looking back at the prose.




\begin{example}[A worked calculation for Scattering, Phase Space, and Cross Sections: Review problems and extensions: checking symmetry factors]
In \(\lambda\phi^4\) theory, the four-point function at first order has one interaction insertion.  The four fields at the vertex can be paired with four external fields in \(4!\) ways, cancelling the \(1/4!\).  The remaining factor is \(-\ii\lambda\), multiplied by four propagators from the vertex point to the external points.  In Scattering, Phase Space, and Cross Sections, section 8, the useful habit is to name the object, the operation, and the assumption before using the compact notation.
\end{example}



\begin{exercise}
Why does the factor \(1/4!\) in \(-\lambda\phi^4/4!\) disappear from the four-point vertex rule?  Use the notation of Scattering, Phase Space, and Cross Sections: Review problems and extensions: checking symmetry factors.
\begin{answer}
There are \(4!\) contractions connecting the four identical fields at the vertex to four external fields.
\end{answer}
\end{exercise}

\begin{exercise}
State what LSZ removes from a connected correlation function.  Use the notation of Scattering, Phase Space, and Cross Sections: Review problems and extensions: checking symmetry factors.
\begin{answer}
It removes the external propagator poles and leaves the on-shell scattering amplitude.
\end{answer}
\end{exercise}



\section{Cumulative worked derivation: from expanding the interaction exponential to identifying loop variables}

This final worked section braids together expanding the interaction exponential, conserving momentum at a vertex, and identifying loop variables for Scattering, Phase Space, and Cross Sections.  The vocabulary of the chapter was S-matrix, rates, phase space, normalization; here those words are used in one continuous calculation rather than in isolated fragments.

Compute the connected four-point function in \(\lambda\phi^4\) theory to first order.  Expand the interaction exponential once.  Contract the four fields at the interaction point with the four external fields.  The \(4!\) possible contractions cancel the \(1/4!\) in the Lagrangian.  The result is
\[
  G_4^{(1)}(x_1,x_2,x_3,x_4)=
  -\ii\lambda\int\dd^4z\,\prod_{r=1}^4\Delta_F(x_r-z).
\]
Fourier transforming this expression gives momentum conservation at the vertex and the vertex factor \(-\ii\lambda\).  In Scattering, Phase Space, and Cross Sections, cumulative derivation, the useful habit is to name the object, the operation, and the assumption before using the compact notation.

For reference, the compact formulas are

\[
  \Delta_F(p)=\frac{\ii}{p^2-m^2+\ii\epsilon}
\]
\[
  e^{\ii\int\Lag_I}=1+\ii\int\Lag_I+\frac{\ii^2}{2}\left(\int\Lag_I\right)^2+\cdots
\]
\[
  \ii\mathcal M=\hbox{sum of connected amputated diagrams}
\]

\begin{exercise}
Reproduce the cumulative derivation for Scattering, Phase Space, and Cross Sections without looking at the prose.  Mark the step where a finite calculation becomes a continuum, operator, or field-theoretic statement.
\begin{answer}
The reconstruction should name the starting object, carry out the differentiating, varying, transforming, or commuting step explicitly, and state the normalization or boundary condition that makes the final formula meaningful.
\end{answer}
\end{exercise}



\section{Chapter synthesis}

This chapter has treated scattering, phase space, and cross sections as a chain of calculations.  The safest summary is not a list of final formulas, but a map from assumptions to equations: finite model, limiting step, boundary condition, normalization, and field-theory use.  If one of those links is missing, the formula should be rederived before it is used in a later chapter.

\begin{exercise}
Write a one-page reconstruction of scattering, phase space, and cross sections.  Include one equation from the finite model, one equation from the continuum or operator form, and one sentence explaining how the two are connected.
\begin{answer}
A strong reconstruction names the basic objects, identifies the central derivative, variation, commutator, or symmetry operation, and states the assumption that allows the finite calculation to become the field-theory formula.
\end{answer}
\end{exercise}
